\section{The part of the argument we can reuse}\label{sec:orbit-criterion}
The local sieves above average a bad event over an orbit, apply a spectral
bound, and add the local losses. This appendix isolates that argument,
allowing the inverse-unit mass needed when several ideal classes occur.
The arithmetic work identifies the orbit, counts the representations,
and leaves a strict margin below one.

The first statement isolates the averaging step. The second lists the
uniform bounds that turn it into an eventual, and potentially effective,
representation theorem.

\begin{proposition}\label{prop:orbit-averaging}
Let a finite group $G$ act on the right on a finite set $Y$ equipped
with a $G$-invariant probability measure $\pi$. Let $K\le G$, and let
$\Omega=y_0K$ be a $K$-orbit of positive measure. On $L^2(Y,\pi)$ put
$f=1_\Omega$ and $\mu=\pi(\Omega)$. Suppose $T$ is self-adjoint,
commutes with the $G$-action, and satisfies $T1=d1$ with $d>0$.
Let $E\subset L^2(Y,\pi)$ be a subspace such that
\[
\|Tv\|\le\Lambda\|v\|\quad(v\in E^\perp),\qquad
f-\mu1\in E^\perp,
\]
where $\Lambda\ge0$. Then
\[
\left|\frac{(Tf)(y_0)}d-\mu\right|
\le\frac\Lambda d(1-\mu).
\]
In particular, the same right side bounds the upper deviation.
\end{proposition}
\begin{proof}
Both $f$ and $Tf$ are right-$K$ invariant, so $Tf$ is constant on
$\Omega$ and
\[
(Tf)(y_0)=\frac{\langle f,Tf\rangle_\pi}{\mu}.
\]
Write $f=\mu1+h$. Then $h\perp1$ and
$\|h\|_\pi^2=\mu(1-\mu)$. Self-adjointness and $T1=d1$ remove the
cross terms, giving
\[
\langle f,Tf\rangle_\pi=d\mu^2+\langle h,Th\rangle_\pi.
\]
The spectral hypothesis bounds the absolute value of the last term by
$\Lambda\mu(1-\mu)$. Divide by $d\mu$.
\end{proof}

Global transitivity is unnecessary: only constancy on the orbit
$\Omega$ enters the averaging identity. In Section~\ref{sec:minus11},
the principal ideal-class fiber has probability $w=3/5$, and the inert
bad orbit has conditional density $\mu_p=2/(p(p-1))$ in that fiber.
Thus its measure in the proposition is $\mu=w\mu_p$.
The orthogonality to norm characters and the separate ramified
calculation are checked there explicitly.

\begin{proposition}\label{prop:arithmetic-criterion}
For each admissible positive integer $n\ge n_0$, let $\mathcal R_n$ be a finite
set of $R(n)>0$ candidate representations. Let $E_n$ be an exceptional
bad subset, and $B_{p,n}$ bad subsets indexed by a fixed set $\mathcal S$
of primes. Every representation outside
$E_n\cup\bigcup_{p\in\mathcal S}B_{p,n}$ is assumed to be good.
Suppose, uniformly in the admissible targets and primes, that
\[
\frac{|E_n|}{R(n)}\le\rho+\eta(n),\qquad
\frac{|B_{p,n}|}{R(n)}\le\mu_p+C_1n^{\delta-1/2},\qquad
\frac{|B_{p,n}|}{R(n)}\le C_2n^\delta p^{-2},
\]
where $C_1,C_2,\rho,\mu_p,\eta(n)\ge0$, $\delta<1/4$,
$\eta(n)\to0$, and
\[
\sum_{p\in\mathcal S}\mu_p\le M<\infty,\qquad \rho+M<1.
\]
Assume also, for every real $P\ge P_0>0$, that
\[
\#\{p\in\mathcal S:p\le P\}\le\alpha P+\beta,\qquad
\sum_{\substack{p\in\mathcal S\\p>P}}p^{-2}
\le\frac\gamma P+\frac\zeta{P^2},
\]
with nonnegative constants. For fixed $t>0$, put
\[
L_1=\alpha tC_1+\frac{\gamma C_2}{t},\qquad
L_2=\beta C_1+\frac{\zeta C_2}{t^2}.
\]
Every admissible $n\ge\max\{n_0,(P_0/t)^4\}$ satisfying
\[
\eta(n)+L_1n^{\delta-1/4}+L_2n^{\delta-1/2}<1-\rho-M
\]
has a good representation. In particular, every sufficiently large
admissible target does.

If the constants are effectively known and $\eta$ has an effective
modulus $N_\eta(u)$, meaning $\eta(n)\le u$ for all admissible
$n\ge N_\eta(u)$, put $\Delta=1-\rho-M$, $a=1/4-\delta$ and
$b=1/2-\delta$. An explicit sufficient cutoff is
\[
\max\left\{n_0,(P_0/t)^4,N_\eta(\Delta/4),
 (4L_1/\Delta)^{1/a},(4L_2/\Delta)^{1/b}\right\}.
\]
\end{proposition}
\begin{proof}
Use the first local bound for $p\le P=t n^{1/4}$ and the second for
$p>P$. The union bound gives
\begin{align*}
\frac{|E_n|+\sum_{p\in\mathcal S}|B_{p,n}|}{R(n)}
&\le\rho+M+\eta(n)+C_1n^{\delta-1/2}(\alpha P+\beta)
+C_2n^\delta\left(\frac\gamma P+\frac\zeta{P^2}\right)\\
&=\rho+M+\eta(n)+L_1n^{\delta-1/4}+L_2n^{\delta-1/2}.
\end{align*}
The tail bound makes the prime sum convergent. If the right side is
less than one, the bad subsets cannot cover $\mathcal R_n$.
Both powers tend to zero because $\delta<1/4$.
At the displayed effective cutoff, the three error terms are each at
most $\Delta/4$, giving a total at most $1-\Delta/4<1$.
\end{proof}
The prime-count hypotheses always admit the crude choices
$\alpha=\gamma=\zeta=1$, $\beta=0$, $P_0=1$, by comparison with integers
and a decreasing integral. Arithmetic residue information improves these
constants. When $\alpha C_1$ and $\gamma C_2$ are positive, the leading
coefficient $L_1$ is minimized at
\[
t=\sqrt{\frac{\gamma C_2}{\alpha C_1}}.
\]
This minimizes that coefficient, not necessarily the complete cutoff.
The mixed Eisenstein application uses $\delta=1/5$, $\rho=1/3$,
$\eta(n)=2/(3\sqrt n)$, $C_1=9/4$, $C_2=15/2$, $M=3/20$ and
$(\alpha,\beta,\gamma,\zeta)=(2/15,17/15,2/15,8/5)$.
Section~\ref{sec:infinite} uses $t=2$ and the sharper direct inequality
there to obtain $10^8$.
